\section{数据审查与建模边界：保留原谱和未知条件}

上一章区分了周期等效量与拟合厚度，本章先核对支撑这些结果的原始光谱，再确定模型使用的观测范围。附件口径见表~\ref{tab:attachment-check}，共同坐标与拟合区段见附录图~\ref{fig:common-coordinate}。
% src:交接/数据档案.json:总览.附件总数;总览.独立晶圆数;总览.每块晶圆入射角度;跨文件关系.同片分组

异常原值遵循“标记、保留、单列影响”。附件二有262个超出反射率物理范围的原始值，集中在低波数端；图\ref{fig:raw-anomaly}和表\ref{tab:attachment-check}保留其位置与幅值。第一次尝试按物理范围裁剪，核对后改为保留原值，因为异常段与主拟合窗口不相交；保留或裁剪这些窗外点，主拟合窗口内的输入均相同。后文据此展开附件体检与参数缺口核对。
% src:交接/数据档案.json:总览.超过百分之百总点数;文件档案.1.工作表.0.字段档案.1.统计.最大值;文件档案.1.工作表.0.反射率检查.超过百分之百区段
% src:交接/实验记录.json:535.现象;535.决定
% src:交接/数据档案.json:未随附件提供的建模输入
